% ============================================================
%  Experiment 9 --- Project Creation in SCADA (IoT)
%  Sourced from the department's i-Factory IIoT Training Manual,
%  Module 3: Project Creation in SCADA.
% ============================================================
\experimentheader{9}{Experiment 9}{Project Creation in SCADA (IoT)}{Instrumentation Lab}

\begin{objectives}
  \begin{enumerate}
    \item To understand the role and capabilities of SCADA software in an industrial IoT system.
    \item To create a working SCADA project, including a Project Node and a SCADA Node.
    \item To configure a communication link (Comport) between the SCADA node and an automation device.
    \item To create and configure IO Tags and Internal Tags for real-time data acquisition.
    \item To configure basic alarms on analog and discrete tags.
  \end{enumerate}
\end{objectives}

\begin{equipment}
  \begin{itemize}
    \item PC/laptop with Advantech WebAccess/SCADA installed (Project Node and SCADA Node)
    \item Siemens S7-1200 PLC (or equivalent), with an existing ladder-logic program and known I/O/DB addressing
    \item Ethernet cable and a configured local network connection between the PC and the PLC
    \item TIA Portal (or equivalent Siemens engineering software) for reference to the PLC's tag addresses
  \end{itemize}
\end{equipment}

\begin{introduction}
  Supervisory Control and Data Acquisition (SCADA) software sits between the plant floor and the people (or systems) that need to monitor and control it. It collects real-time data from automation devices --- PLCs, controllers, direct digital control (DDC) systems, and distributed control systems (DCS) --- and presents that data through live displays, trends, alarms, and reports.
  \vspace*{1.5em}

  Advantech WebAccess/SCADA, used in this experiment, is a 100\% web-based SCADA package: projects are developed and accessed through a standard web browser, rather than a dedicated desktop client. This lowers the barrier to remote monitoring and multi-site deployment, and it integrates directly with IoT dashboards and cloud platforms via MQTT.
  \vspace*{1.5em}

  A WebAccess/SCADA system is built from three components:
  \begin{itemize}
    \item \textbf{Project Node} --- the development platform and web server. All project configuration, graphics, and database files are stored here, and every client connects to it first.
    \item \textbf{SCADA Node} --- runs the real-time Kernel (\texttt{datacore.exe}), which uses communication drivers to talk to automation hardware (PLCs, I/O modules, controllers) and supplies live data to clients.
    \item \textbf{ViewDAQ Client} --- the browser-based interface used to monitor and control the SCADA node: viewing live graphics, historical trends, and alarms, and adjusting set points.
  \end{itemize}
  In a small setup (as in this lab), the Project Node and SCADA Node commonly run on the same machine.
  \vspace*{1.5em}

  Data inside WebAccess/SCADA is represented as \textbf{Tags}. There are two broad categories:
  \begin{itemize}
    \item \textbf{IO Tags} --- real-time measurements and outputs read from, or written to, an automation device. IO Tags can be Analog (floating-point), Discrete (an integer 0--7, commonly used as a 0/1 digital value), or Text (an ASCII string up to 72 characters).
    \item \textbf{Internal Tags} --- values that exist inside the SCADA software itself rather than being read from a device: \textbf{Constant} tags (a user-defined fixed or operator-editable value), \textbf{Calculation} tags (a mathematical/logical expression over up to 20 other tags), and \textbf{Accumulation} tags (an integration/totaliser over another tag, e.g.\ totalising flow into volume).
  \end{itemize}
  Both tag categories share the same alarm, security, logging, and reporting features once created.
\end{introduction}

\needspace{6cm}
\begin{boxedtable}
  \captionof{table}{Example Siemens S7-1200 address mapping to WebAccess tag addressing.}
  \label{tab:exp9scada-addrmap}
  {\renewcommand{\arraystretch}{1.4}
    \centering
    \begin{tabular}{llcl}
      \toprule
      Siemens S7-1200 address & Data type & Length (bits) & WebAccess address    \\
      \midrule
      \%I0.4                 & Bool      & 1              & IX000 (start bit 4)  \\
      \%Q0.3                 & Bool      & 1              & QX000 (start bit 3)  \\
      \%M0.4                 & Bool      & 1              & MX000 (start bit 4)  \\
      \%MW0                  & UInt      & 16             & MW000 (start bit 0)  \\
      \%IW6.4                & UInt      & 16             & IW006 (start bit 4)  \\
      \%DB1.DBX2.0            & Bool      & 1              & DBX1.2 (start bit 0) \\
      \%DB3.DBW10             & UInt      & 16             & DBW3,10               \\
      \%DB3.DBD14             & Real      & 64             & DBD3,14               \\
      \bottomrule
    \end{tabular}
  }
  \vspace{4pt}

  \small\textit{Note: for tags addressed inside a Data Block (DB), the offset must be enabled in the block's properties in TIA Portal, since WebAccess addresses DB tags by block number plus offset.}
\end{boxedtable}

\begin{procedure}

  \textbf{Activity 1 --- Create a new SCADA project}\vspace{2pt}

  \textit{Objective: create a Project Node and SCADA Node in WebAccess/SCADA.}

  \begin{enumerate}
    \item Open a web browser and navigate to \texttt{127.0.0.1} (or \texttt{localhost}) to reach the WebAccess/SCADA landing page.
    \item Create a new project: enter a \textbf{Project Name} and \textbf{Project Description}, and set the \textbf{Project Node IP Address} (use \texttt{127.0.0.1} for a local machine, or the PC's hostname).
    \item Leave the \textbf{Primary TCP Port} (default \texttt{4592}) and \textbf{HTTP Port} (default \texttt{80}) at their defaults unless your network/firewall requires otherwise, then click \textbf{Submit}.
    \item Select the Project Node in the tree structure, click \textbf{+}, and choose \textbf{New Node} to create the SCADA Node. Fill in the Primary Node Name, Primary/Secondary TCP Port, Node Timeout, and SCADA Node IP Address.
    \item Start the Kernel on the SCADA Node (\textbf{Start Kernel} from the SCADA Node menu) and confirm it is running by hovering over the WebAccess/SCADA icon, which should report the running project and node name.
  \end{enumerate}
  \vspace*{1.5em}

  \textbf{Activity 2 --- Configure a communication port and add a device}\vspace{2pt}

  \textit{Objective: create the physical/virtual link (Comport) between the SCADA node and the PLC, and add the PLC as a device on that port.}

  \begin{enumerate}
    \item Under the SCADA Node, add a new \textbf{Comport}. Select the interface type appropriate to your connection:
          \begin{itemize}
            \item \textbf{Serial} --- an RS232/RS422/RS485 port (COM1, COM2, \ldots) on the SCADA node PC.
            \item \textbf{TCP/IP} --- the most common choice for PLCs and IoT devices; WebAccess/SCADA searches all TCP/IP network cards regardless of the assigned Comport number.
            \item \textbf{OPC} --- via an OPC Server, for virtual/software-based links.
          \end{itemize}
    \item Configure the port: Port Number, Scan Time (e.g.\ \texttt{1000\,ms}), and Auto-recover Time.
    \item Under the configured Comport, add a new \textbf{Device}, selecting the correct Device Type/Driver for your hardware (e.g.\ Siemens S7-1200, Modbus TCP, Allen-Bradley).
    \item Enter the device's communication parameters (IP address, station ID, or baud rate for serial links) and save.
  \end{enumerate}
  \vspace*{1.5em}

  \textbf{Activity 3 --- Create IO Tags}\vspace{2pt}

  \textit{Objective: create IO Tags that map to real signals in the PLC program.}

  \begin{enumerate}
    \item Identify at least four signals in your PLC program to bring into SCADA --- for example, a Start pushbutton, a Stop pushbutton, an Emergency-stop input, and one analog or word-type value (e.g.\ a counter or measured value).
    \item For each signal, determine its native PLC address and convert it to the equivalent WebAccess address using Table~\ref{tab:exp9scada-addrmap} as a reference (e.g.\ \texttt{\%DB1.DBX0.1} for a Start pushbutton, \texttt{\%DB1.DBX0.3} for Stop, \texttt{\%DB1.DBX0.0} for E-stop).
    \item Under the Device, click the tag icon and \textbf{Insert} a new tag for each signal. Fill in the Tag Type, Tag Name, Address, Start Bit, Conversion Code, and Length fields to match your PLC addressing.
    \item Save your tags, then verify each one updates correctly by forcing the corresponding input/output in the PLC and observing the tag's live value in WebAccess.
  \end{enumerate}
  \vspace*{1.5em}

  \textbf{Activity 4 --- Create Internal Tags}\vspace{2pt}

  \textit{Objective: create at least one Constant Tag and one Calculation Tag.}

  \begin{enumerate}
    \item Under \textbf{Const Point} on the SCADA Node, insert a new Constant Tag (analog type, \texttt{ConAna}, or discrete type, \texttt{ConDis}) --- for example, a user-adjustable setpoint value.
    \item Under \textbf{Calc Point} on the SCADA Node, insert a new Calculation Tag. Choose an analog or discrete result type, then write a mathematical/logical expression referencing one or more of your IO Tags (tags are referenced in the formula by their single-letter designations; the expression is limited to 80 characters).
    \item Save and verify that the Calculation Tag's value updates correctly as its input tag(s) change.
  \end{enumerate}
  \vspace*{1.5em}

  \textbf{Activity 5 --- Configure a basic alarm}\vspace{2pt}

  \textit{Objective: enable alarming on at least one analog tag and one discrete tag.}

  \begin{enumerate}
    \item Open the configuration for one of your Analog IO Tags and enable a non-zero alarm priority. Analog tags support High-High, High, Low, Low-Low, Deviation, and Rate-of-Change alarms --- choose whichever is appropriate to the signal you selected.
    \item Open the configuration for one of your Discrete IO Tags and enable state alarming (e.g.\ alarm when the E-stop input becomes active).
    \item Trigger each alarm condition on the PLC side and confirm that the alarm appears correctly in WebAccess.
  \end{enumerate}

\end{procedure}

\begin{reportq}
  \begin{enumerate}
    \item List the Project Node and SCADA Node settings you used (IP address, ports, node names), and explain the role each plays in the system.
    \item Present your Comport and Device configuration, including which interface type (Serial/TCP-IP/OPC) you selected and why it was appropriate for your hardware.
    \item Tabulate the IO Tags you created: PLC address, WebAccess address, tag type, and a description of what each represents.
    \item Present your Constant Tag and Calculation Tag configurations, including the expression used for the Calculation Tag and what it computes.
    \item Describe the alarm(s) you configured and demonstrate (with a screenshot or description of your test) that each triggers correctly.
    \item Discuss one advantage of WebAccess/SCADA's 100\% web-based architecture compared to a traditional desktop-client SCADA system.
  \end{enumerate}
\end{reportq}

\begin{labsection}{References}{}
\begin{itemize}
  \item MTA SkillLab Sdn Bhd, \textit{i-Factory Training Manual --- IIoT System Development Level 1}, Module 3: Project Creation in SCADA.
  \item Advantech WebAccess/SCADA product documentation: \url{https://www.advantech.com/en/support/details/installation?id=1-MS9MJV}
\end{itemize}
\end{labsection}
